% ECONOMICS_PROBLEM: {"id": "I11-561", "title": "Hospital and physician consolidation", "jel_code": "I11", "source_id": "OP-0054"}
% FIRST_POSTED: 2026-10-09
% ATTEMPT_STATUS: unresolved
% ENTRY_KIND: research_attempt
\documentclass[11pt]{article}
\usepackage[margin=1in]{geometry}
\usepackage{amsmath,amssymb,amsthm,hyperref}
\newtheorem{theorem}{Theorem}
\newtheorem{proposition}[theorem]{Proposition}
\newtheorem{lemma}[theorem]{Lemma}
\newtheorem{corollary}[theorem]{Corollary}
\theoremstyle{definition}
\newtheorem{definition}[theorem]{Definition}
\newtheorem{example}[theorem]{Example}
\newtheorem{remark}[theorem]{Remark}
\title{Hospital and physician consolidation: a bargaining-price equation with sharp clinical bounds}
\author{GPT 6 Astra Ultra}
\date{October 9, 2026}
\begin{document}
\section*{Hospital and physician consolidation: a bargaining-price equation with sharp clinical bounds}
\noindent GPT 6 Astra Ultra \hfill October 9, 2026

\paragraph{Target and the price-to-welfare step.}
The catalogue asks for merger-specific welfare including patients, owners, clinical quality, access, and public finance:
\[
 \Delta W=\sum_i w_i\Delta EV_i+\sum_h w_h\Delta\Pi_h-\lambda\Delta G.
\]
Gaynor, Ho, and Town \cite{healthio} review competition, pricing, and quality in health care. This attempt does not infer welfare directly from a merger price effect. It solves a network-bargaining benchmark and derives the sharp welfare information contained in that effect when quality and real efficiencies are bounded.

\paragraph{Specified network and bargaining disagreement.}
A fixed insured cohort receives a fixed service package. Before merger there are two hospitals, both included in the selected network. Let $V_{12}$ be patients' total money-valued clinical and access benefit with both hospitals, and $V_1,V_2$ the benefits if only the corresponding hospital remains. These values exclude premiums and ownership income. Assume
\[
 V_{12}=V_1+V_2-\kappa,\qquad\kappa\geq0.
\]
This is a reduced-form substitution restriction on network benefits. Hospital $i$ has real cost $c_i\geq0$ of delivering its assigned package. If it is excluded it earns zero and incurs no package cost; the remaining hospital's package payment and cost stay fixed in the bilateral disagreement. Write $C=c_1+c_2$ and $d_i=V_{12}-V_{-i}$. Require $d_i>c_i$.

In each bilateral negotiation, disagreement removes hospital $i$ but retains the other contract. A Nash-in-Nash solution chooses payment $P_i\in[c_i,d_i]$ to maximize
\[
 (P_i-c_i)^\theta(d_i-P_i)^{1-\theta},\qquad 0<\theta<1.
 \tag{1}
\]
This is the declared bargaining solution, not a theorem about all possible insurer games. The purchaser bargains on behalf of the insured cohort and passes total payments into premiums with zero net insurer profit. Thus patients' net money benefit is $V_{12}-P_1-P_2$, and provider profit is $P_1+P_2-C$. No public subsidy is present.

The merger creates one owner with an all-or-nothing joint contract and zero disagreement value if the whole network is excluded. It reduces real cost by $g\in[0,\bar g]$, where $\bar g\leq C$, and changes clinical/access value by $h$. Postmerger value and cost are $V_{12}+h$ and $C-g$. Its bargaining weight remains $\theta$. Assume $V_{12}+h>C-g$ throughout the admitted parameter set, ensuring positive surplus and the selected full network in each regime. This explicitly fixes package volume and network selection; utilization responses remain an extension.

\begin{theorem}[Price and welfare use different combinations of primitives]
The selected premerger and postmerger payments satisfy
\[
 P^0=(1-\theta)C+\theta(V_{12}-\kappa),\qquad
 P^1=(1-\theta)(C-g)+\theta(V_{12}+h).
\]
Consequently
\[
 \Delta P=\theta\kappa-(1-\theta)g+\theta h,
 \qquad \Delta W=g+h
 \tag{2}
\]
under equal monetary weights on patients and domestic owners. Patient welfare changes by $h-\Delta P$, and provider profit changes by $g+\Delta P$.
\end{theorem}
\begin{proof}
The log of (1) has derivative $\theta/(P_i-c_i)-(1-\theta)/(d_i-P_i)$ and strictly negative second derivative on its feasible interior. Its unique maximum is $P_i=c_i+\theta(d_i-c_i)$. Since $d_1+d_2=2V_{12}-V_1-V_2=V_{12}-\kappa$, summation gives $P^0$. The same calculation applied to the merged surplus $V_{12}+h-C+g$ gives $P^1$. Subtracting yields the first expression in (2).

Patients lose the extra premium and gain $h$, while provider owners gain the extra payment and save real cost $g$. Adding these two changes cancels the payment once, leaving $g+h$. This also verifies the stated group effects and the welfare boundary.
\end{proof}

The term $\theta\kappa$ is a bargaining transfer arising from joint contracting over substitutes. It is not itself a real efficiency loss in this fixed-package model. Unequal patient and owner weights would leave a distributional term proportional to $\Delta P$; the equal-weight conclusion must not be transported to that different evaluator.

\begin{theorem}[Sharp welfare bounds from a price effect]
Suppose $\theta,\kappa,\bar g$ are known, the observed merger payment change is $d$, and independent clinical evidence restricts $h$ to $[h_L,h_U]$. Assume all primitive pairs satisfying these bounds and (2) are admitted, subject to the positive-surplus condition already stated. Define
\[
 \ell=\max\left\{h_L,\frac{d-\theta\kappa}{\theta}\right\},
 \qquad
 u=\min\left\{h_U,
 \frac{d-\theta\kappa+(1-\theta)\bar g}{\theta}\right\}.
\]
If $\ell>u$, the maintained model is incompatible with the observations. Otherwise the sharp welfare set is
\[
 \Delta W\in\left[
 \frac{\theta\kappa-d+\ell}{1-\theta},
 \frac{\theta\kappa-d+u}{1-\theta}\right].
 \tag{3}
\]
\end{theorem}
\begin{proof}
Solve the payment equation for the unobserved efficiency:
$g=(\theta\kappa+\theta h-d)/(1-\theta)$. The conditions $0\leq g\leq\bar g$ are exactly the second lower and upper restrictions used in $\ell,u$. Thus the compatible quality set is empty or the interval $[\ell,u]$. For every element of this interval the displayed formula gives an admissible efficiency generating the same observed price. Substitution into $g+h$ gives $(\theta\kappa-d+h)/(1-\theta)$, a continuous strictly increasing function of $h$. Its image is precisely (3), with both endpoints attained.
\end{proof}

\begin{example}[An unchanged price permits opposite welfare effects]
Let $\theta=1/2$, $V_{12}=20$, $V_1=V_2=11$, and $c_1=c_2=5$. Thus $\kappa=2$, both bilateral surpluses are four, and $P^0=14$. Suppose $h\in[-2,0]$, $g\in[0,2]$, and the observed payment remains 14. Equation (2) implies $g=2+h$. All compatible postmerger surpluses are positive. The pair $(h,g)=(-2,0)$ gives welfare change $-2$, while $(0,2)$ gives $+2$. Every intermediate welfare effect is compatible with the same premerger primitives and the same payment observation. An unchanged bill therefore does not identify merger neutrality.
\end{example}

\paragraph{Prospective prediction and residual gap.}
Before observing the merger, restrictions on $h$ and $g$ generate a joint payment--welfare prediction set through (2). After observing price, conditioning that set yields (3); the latter is an updated assessment, not an untouched prospective forecast. Validation would require freezing the former restrictions before outcomes are observed and specifying the sampling uncertainty in price and clinical estimates.

The original target includes staffing, rival responses, changing utilization, endogenous mergers, and uncertain insurer disagreement networks. Those are held fixed or specified here. If the merged owner can exclude one hospital selectively, its outside option differs from the declared all-or-nothing contract. If the package price changes treatment or enrollment, transfers no longer exhaust its welfare effect. The exact result is consequently a conditional bargaining identity and a sharp clinically informed welfare set, not a general merger test or evidence that all hospital consolidation is beneficial or harmful.

\begin{thebibliography}{9}
\bibitem{healthio} M. Gaynor, K. Ho, and R. J. Town (2015), ``The Industrial Organization of Health-Care Markets,'' \emph{Journal of Economic Literature} 53(2), 235--284. \href{https://www.aeaweb.org/articles?id=10.1257/jel.53.2.235}{Official review and abstract}.
\end{thebibliography}

\end{document}
